% Options for packages loaded elsewhere
\PassOptionsToPackage{unicode}{hyperref}
\PassOptionsToPackage{hyphens}{url}
\documentclass[
  ignorenonframetext,
]{beamer}
\newif\ifbibliography
\usepackage{pgfpages}
\setbeamertemplate{caption}[numbered]
\setbeamertemplate{caption label separator}{: }
\setbeamercolor{caption name}{fg=normal text.fg}
\beamertemplatenavigationsymbolsempty
% remove section numbering
\setbeamertemplate{part page}{
  \centering
  \begin{beamercolorbox}[sep=16pt,center]{part title}
    \usebeamerfont{part title}\insertpart\par
  \end{beamercolorbox}
}
\setbeamertemplate{section page}{
  \centering
  \begin{beamercolorbox}[sep=12pt,center]{section title}
    \usebeamerfont{section title}\insertsection\par
  \end{beamercolorbox}
}
\setbeamertemplate{subsection page}{
  \centering
  \begin{beamercolorbox}[sep=8pt,center]{subsection title}
    \usebeamerfont{subsection title}\insertsubsection\par
  \end{beamercolorbox}
}
% Prevent slide breaks in the middle of a paragraph
\widowpenalties 1 10000
\raggedbottom
\AtBeginPart{
  \frame{\partpage}
}
\AtBeginSection{
  \ifbibliography
  \else
    \frame{\sectionpage}
  \fi
}
\AtBeginSubsection{
  \frame{\subsectionpage}
}
\usepackage{iftex}
\ifPDFTeX
  \usepackage[T1]{fontenc}
  \usepackage[utf8]{inputenc}
  \usepackage{textcomp} % provide euro and other symbols
\else % if luatex or xetex
  \usepackage{unicode-math} % this also loads fontspec
  \defaultfontfeatures{Scale=MatchLowercase}
  \defaultfontfeatures[\rmfamily]{Ligatures=TeX,Scale=1}
\fi
\usepackage{lmodern}
\ifPDFTeX\else
  % xetex/luatex font selection
\fi
% Use upquote if available, for straight quotes in verbatim environments
\IfFileExists{upquote.sty}{\usepackage{upquote}}{}
\IfFileExists{microtype.sty}{% use microtype if available
  \usepackage[]{microtype}
  \UseMicrotypeSet[protrusion]{basicmath} % disable protrusion for tt fonts
}{}
\makeatletter
\@ifundefined{KOMAClassName}{% if non-KOMA class
  \IfFileExists{parskip.sty}{%
    \usepackage{parskip}
  }{% else
    \setlength{\parindent}{0pt}
    \setlength{\parskip}{6pt plus 2pt minus 1pt}}
}{% if KOMA class
  \KOMAoptions{parskip=half}}
\makeatother
\usepackage{color}
\usepackage{fancyvrb}
\newcommand{\VerbBar}{|}
\newcommand{\VERB}{\Verb[commandchars=\\\{\}]}
\DefineVerbatimEnvironment{Highlighting}{Verbatim}{commandchars=\\\{\}}
% Add ',fontsize=\small' for more characters per line
\newenvironment{Shaded}{}{}
\newcommand{\AlertTok}[1]{\textcolor[rgb]{1.00,0.00,0.00}{\textbf{#1}}}
\newcommand{\AnnotationTok}[1]{\textcolor[rgb]{0.38,0.63,0.69}{\textbf{\textit{#1}}}}
\newcommand{\AttributeTok}[1]{\textcolor[rgb]{0.49,0.56,0.16}{#1}}
\newcommand{\BaseNTok}[1]{\textcolor[rgb]{0.25,0.63,0.44}{#1}}
\newcommand{\BuiltInTok}[1]{\textcolor[rgb]{0.00,0.50,0.00}{#1}}
\newcommand{\CharTok}[1]{\textcolor[rgb]{0.25,0.44,0.63}{#1}}
\newcommand{\CommentTok}[1]{\textcolor[rgb]{0.38,0.63,0.69}{\textit{#1}}}
\newcommand{\CommentVarTok}[1]{\textcolor[rgb]{0.38,0.63,0.69}{\textbf{\textit{#1}}}}
\newcommand{\ConstantTok}[1]{\textcolor[rgb]{0.53,0.00,0.00}{#1}}
\newcommand{\ControlFlowTok}[1]{\textcolor[rgb]{0.00,0.44,0.13}{\textbf{#1}}}
\newcommand{\DataTypeTok}[1]{\textcolor[rgb]{0.56,0.13,0.00}{#1}}
\newcommand{\DecValTok}[1]{\textcolor[rgb]{0.25,0.63,0.44}{#1}}
\newcommand{\DocumentationTok}[1]{\textcolor[rgb]{0.73,0.13,0.13}{\textit{#1}}}
\newcommand{\ErrorTok}[1]{\textcolor[rgb]{1.00,0.00,0.00}{\textbf{#1}}}
\newcommand{\ExtensionTok}[1]{#1}
\newcommand{\FloatTok}[1]{\textcolor[rgb]{0.25,0.63,0.44}{#1}}
\newcommand{\FunctionTok}[1]{\textcolor[rgb]{0.02,0.16,0.49}{#1}}
\newcommand{\ImportTok}[1]{\textcolor[rgb]{0.00,0.50,0.00}{\textbf{#1}}}
\newcommand{\InformationTok}[1]{\textcolor[rgb]{0.38,0.63,0.69}{\textbf{\textit{#1}}}}
\newcommand{\KeywordTok}[1]{\textcolor[rgb]{0.00,0.44,0.13}{\textbf{#1}}}
\newcommand{\NormalTok}[1]{#1}
\newcommand{\OperatorTok}[1]{\textcolor[rgb]{0.40,0.40,0.40}{#1}}
\newcommand{\OtherTok}[1]{\textcolor[rgb]{0.00,0.44,0.13}{#1}}
\newcommand{\PreprocessorTok}[1]{\textcolor[rgb]{0.74,0.48,0.00}{#1}}
\newcommand{\RegionMarkerTok}[1]{#1}
\newcommand{\SpecialCharTok}[1]{\textcolor[rgb]{0.25,0.44,0.63}{#1}}
\newcommand{\SpecialStringTok}[1]{\textcolor[rgb]{0.73,0.40,0.53}{#1}}
\newcommand{\StringTok}[1]{\textcolor[rgb]{0.25,0.44,0.63}{#1}}
\newcommand{\VariableTok}[1]{\textcolor[rgb]{0.10,0.09,0.49}{#1}}
\newcommand{\VerbatimStringTok}[1]{\textcolor[rgb]{0.25,0.44,0.63}{#1}}
\newcommand{\WarningTok}[1]{\textcolor[rgb]{0.38,0.63,0.69}{\textbf{\textit{#1}}}}
\setlength{\emergencystretch}{3em} % prevent overfull lines
\providecommand{\tightlist}{%
  \setlength{\itemsep}{0pt}\setlength{\parskip}{0pt}}
% Slide layout defaults for warning-clean scientific decks.
\usepackage{etoolbox}
\IfFileExists{xurl.sty}{\usepackage{xurl}}{}
\IfFileExists{seqsplit.sty}{\usepackage{seqsplit}}{\newcommand{\seqsplit}[1]{#1}}
\protected\def\breakseq#1{\seqsplit{#1}}
\protected\def\breaktt#1{\begingroup\ttfamily\seqsplit{#1}\endgroup}
\setlength{\emergencystretch}{6em}
\tolerance=5000
\hbadness=10000
\hfuzz=1pt
\setlength{\tabcolsep}{2pt}
\AtBeginEnvironment{longtable}{\tiny\renewcommand{\arraystretch}{0.86}\setlength{\tabcolsep}{1pt}}
\AtBeginEnvironment{tabular}{\tiny\renewcommand{\arraystretch}{0.86}\setlength{\tabcolsep}{1pt}}

% Natbib and cross-reference fallbacks — slides don't load natbib
% or cleveref, but combined-PDF manuscript prose may emit these
% commands. The fallback renders citations as a bracketed key list
% and cross-references as detokenized labels so slides don't fail on
% undefined control sequences or raw underscores. \providecommand is
% a no-op if packages load later.
\providecommand{\citep}[1]{[#1]}
\providecommand{\citet}[1]{#1}
\providecommand{\citealp}[1]{#1}
\providecommand{\citeauthor}[1]{#1}
\providecommand{\citeyear}[1]{#1}
\providecommand{\cref}[1]{\texttt{\detokenize{#1}}}
\providecommand{\Cref}[1]{\texttt{\detokenize{#1}}}
\usepackage{bookmark}
\IfFileExists{xurl.sty}{\usepackage{xurl}}{} % add URL line breaks if available
\urlstyle{same}
% fallback for those not using the hyperref driver hyperxmp:
\makeatletter
\@ifundefined{xmpquote}{\newcommand{\xmpquote}[1]{#1}}{}
\makeatother
\hypersetup{
  hidelinks,
  pdfcreator={LaTeX via pandoc}}

\author{\texorpdfstring{}{}}
\date{}

\begin{document}

\begin{frame}[fragile,allowframebreaks]{Methods}
\protect\phantomsection\label{methods}
\begin{block}{The A→E generation spine}
\protect\phantomsection\label{the-ae-generation-spine}
Generation proceeds through five stages, each implemented as a pure
function in its own module (\texttt{src/grammar.py},
\texttt{src/expand.py}, \breaktt{src/materialize.py},
\texttt{src/verify.py}, \texttt{src/sealing.py}). No stage depends on
interactive state or network access; every stage takes an immutable
input and returns an immutable (or file-system-materialized) output.

\begin{Shaded}
\begin{Highlighting}[]
\NormalTok{flowchart LR}
\NormalTok{    A[Load Grammar] {-}{-}\textgreater{} B[Expand Spec]}
\NormalTok{    B {-}{-}\textgreater{} C[Materialize Child]}
\NormalTok{    C {-}{-}\textgreater{} D[Verify Integrity]}
\NormalTok{    D {-}{-}\textgreater{} E[Seal + QR]}
\end{Highlighting}
\end{Shaded}
\end{block}

\begin{block}{Stage A --- Grammar loading and validation}
\protect\phantomsection\label{stage-a-grammar-loading-and-validation}
\breaktt{load\_grammar(project\_root)} reads
\breaktt{manuscript/config.yaml}, extracts the \texttt{autopoiesis:}
block, and hands it to \texttt{parse\_grammar()}. Parsing enforces four
invariants before a \texttt{Grammar} object can exist at all:

\begin{enumerate}
\tightlist
\item
  \textbf{The seed must be an integer.} \texttt{parse\_grammar} reads
  \breaktt{block.get("seed",\ 42)} and raises \texttt{GrammarError} if
  the value is not an \texttt{int} --- a stray string or float seed in
  \texttt{config.yaml} fails loudly at load time rather than silently
  coercing.
\item
  \textbf{At least one slot must be defined.} An empty or missing
  \texttt{slots:} list raises
  \breaktt{GrammarError("Grammar\ must\ define\ at\ least\ one\ slot")}.
\item
  \textbf{Every slot must have a non-empty name and ≥1 option.} These
  checks live in \breaktt{GrammarSlot.\_\_post\_init\_\_}, so they fire
  the instant a \texttt{GrammarSlot} is constructed --- a malformed
  entry cannot survive to become part of a \texttt{Grammar}.
\item
  \textbf{No slot may contain duplicate options.}
  \breaktt{GrammarSlot.\_\_post\_init\_\_} also computes
  \texttt{\{o\ for\ o\ in\ self.options\ if\ self.options.count(o)\ \textgreater{}\ 1\}}
  and raises if that set is non-empty, so two identical options in one
  slot (e.g.~\texttt{{[}optimization,\ optimization{]}}) are rejected
  rather than silently collapsing the product space.
\end{enumerate}

\texttt{parse\_grammar} additionally validates every entry in
\texttt{deps:} against \texttt{VENDORABLE\_DEPS} --- the fixed tuple
\breaktt{(logging,\ glossary\_gen,\ figure\_manager,\ manuscript\_injection,\ steganography)}
--- raising \texttt{GrammarError} on any unknown dependency name. This
project's own grammar (see \breaktt{manuscript/config.yaml}) currently
declares \texttt{deps:\ {[}{]}}, so the deps-vendoring path exercised in
\texttt{materialize.py} is present in the code and covered by
\breaktt{test\_deps\_vendoring.py}, but not active for the manuscript's
own default render.

A successfully constructed \texttt{Grammar} is a frozen dataclass
carrying \texttt{seed}, the tuple of \texttt{GrammarSlot} objects, the
tuple of \texttt{deps}, and an optional \texttt{source\_path} (excluded
from equality/hash comparison so two grammars loaded from different
files but with identical content still compare equal).
\end{block}

\begin{block}{Reserved slots vs.~effective slots}
\protect\phantomsection\label{reserved-slots-vs.-effective-slots}
Not every slot in the grammar contributes to what this paper calls the
\emph{effective product size}. \texttt{RESERVED\_SLOTS} is a fixed tuple
--- \texttt{figure\_profile}, \texttt{qr\_profile},
\breaktt{integrity\_profile} --- naming slots that control
\textbf{presentation and provenance mechanics} (how many figures render,
whether a QR seal is embedded, which hash scheme secures the tree hash)
rather than \textbf{domain content} (which primitive kernel, which
analytical track, which manuscript section set). \texttt{Grammar}
exposes both views as properties:

\begin{itemize}
\tightlist
\item
  \texttt{Grammar.slots} --- every slot as declared in
  \texttt{config.yaml}.
\item
  \breaktt{Grammar.reserved\_slots} --- the subset whose \texttt{name}
  appears in \texttt{RESERVED\_SLOTS}.
\item
  \breaktt{Grammar.effective\_slots} --- the complementary subset,
  i.e.~every slot \emph{not} in \texttt{RESERVED\_SLOTS}.
\item
  \breaktt{Grammar.product\_size} --- the raw cross product over
  \emph{all} slots (\breaktt{n\ *=\ len(s.options)} for every slot,
  reserved or not).
\item
  \breaktt{Grammar.effective\_product\_size} --- the cross product
  restricted to \texttt{effective\_slots} only.
\end{itemize}

The distinction matters for honest reporting: a grammar can nominally
claim a large product space by adding presentation-only slots, while the
number of \emph{substantively distinct} generated projects --- different
kernel domain, different analytical track, different section layout ---
is the smaller effective figure. Both \texttt{360} (nominal) and
\texttt{45} (effective) are reported in this manuscript side-by-side
rather than only the larger, more impressive-looking number --- this is
the paper's concrete instance of the ``hard to vary'' honesty discipline
the project holds itself to (see the Honesty Contract, §4). \texttt{3}
reserved slots
(\breaktt{figure\_profile,\ qr\_profile,\ integrity\_profile}) are
excluded from the effective figure; \texttt{SYNTAX.md} documents the
same slot table and the \emph{inflation factor} (nominal ÷ effective) as
a first-class, honestly-reported quantity rather than an implementation
detail.

\breaktt{force\_domain(grammar,\ domain)} provides a targeted override:
it returns a new \texttt{Grammar} with the \breaktt{primitive\_domain}
slot's options collapsed to a single forced value (raising
\texttt{GrammarError} if \texttt{domain} is not one of the five
\texttt{KNOWN\_DOMAINS}), leaving every other slot --- reserved or
effective --- untouched. This is how the pipeline can materialize ``one
child per domain'' deterministically without re-deriving the whole
grammar.
\end{block}

\begin{block}{Stage B --- Deterministic spec expansion}
\protect\phantomsection\label{stage-b-deterministic-spec-expansion}
\breaktt{expand(grammar,\ seed=None)} walks \texttt{grammar.slots} in
declaration order and, for each slot, computes an index into that slot's
options via \texttt{\_digest\_index}:

\begin{Shaded}
\begin{Highlighting}[]
\KeywordTok{def}\NormalTok{ \_digest\_index(seed, slot\_name, ordinal, options):}
\NormalTok{    key }\OperatorTok{=} \SpecialStringTok{f"}\SpecialCharTok{\{}\NormalTok{seed}\SpecialCharTok{\}}\CharTok{\textbackslash{}x1f}\SpecialCharTok{\{}\NormalTok{slot\_name}\SpecialCharTok{\}}\CharTok{\textbackslash{}x1f}\SpecialCharTok{\{}\NormalTok{ordinal}\SpecialCharTok{\}}\CharTok{\textbackslash{}x1f}\SpecialCharTok{\{}\StringTok{\textquotesingle{},\textquotesingle{}}\SpecialCharTok{.}\NormalTok{join(options)}\SpecialCharTok{\}}\SpecialStringTok{"}
\NormalTok{    digest }\OperatorTok{=}\NormalTok{ hashlib.sha256(key.encode()).digest()}
\NormalTok{    value }\OperatorTok{=} \BuiltInTok{int}\NormalTok{.from\_bytes(digest[:}\DecValTok{8}\NormalTok{], }\StringTok{"big"}\NormalTok{)}
    \ControlFlowTok{return}\NormalTok{ value }\OperatorTok{\%} \BuiltInTok{len}\NormalTok{(options)}
\end{Highlighting}
\end{Shaded}

Four inputs --- the seed, the slot's own name, its ordinal position in
the slot list, and the full joined option string --- are concatenated
with an ASCII unit-separator (\texttt{\textbackslash{}x1f}) between
fields and hashed with SHA-256. The first eight bytes of the digest are
read as a big-endian integer and reduced modulo the option count. This
construction has three consequences load- bearing for the rest of the
pipeline:

\begin{itemize}
\tightlist
\item
  \textbf{No shared PRNG state.} Each slot's selection is an independent
  hash of its own key, not a draw from a stateful random-number
  generator advanced slot-by-slot. Reordering unrelated slots elsewhere
  in the list does not perturb a given slot's selection unless that
  slot's own \texttt{ordinal} changes.
\item
  \textbf{Full avalanche on any input change.} Because the selection is
  a SHA-256 digest, changing the seed, renaming a slot, or
  adding/removing a single option string anywhere in that slot's option
  tuple changes the digest --- and therefore, with high probability, the
  selected index --- for that slot.
\item
  \textbf{Reproducibility without stored randomness.} Nothing about the
  selection is persisted except the seed and the grammar itself;
  re-running \texttt{expand} against the same grammar and seed
  recomputes the identical digest and therefore the identical selection,
  with no cached or serialized RNG state to keep in sync.
\end{itemize}

The result is a frozen \texttt{Spec} dataclass: \texttt{schema\_version}
(the fixed string \breaktt{"autopoiesis/spec/1"}), the \texttt{seed}
actually used (the explicit override if supplied, otherwise
\texttt{grammar.seed}), the \texttt{grammar\_hash} this spec was
expanded against, an ordered tuple of
\breaktt{(slot\_name,\ chosen\_value)} \texttt{selections}, the
\texttt{deps} inherited unchanged from the grammar, and the resolved
\breaktt{primitive\_domain} (read out of the selections if a
\breaktt{primitive\_domain} slot exists, else defaulted to
\texttt{KNOWN\_DOMAINS{[}0{]}}). A \texttt{Spec} additionally exposes
\texttt{spec\_hash}, computed by serializing \texttt{to\_dict()} to
canonical (sorted-key, compact-separator) JSON and taking the first 16
hex characters of its SHA-256 --- the same truncation convention used
for \texttt{grammar\_hash}.

Two auxiliary functions extend \texttt{expand} to families of children
rather than one: \breaktt{derive\_seed(base\_seed,\ index)} hashes
\texttt{f"\{base\_seed\}\textbackslash{}x1f\{index\}"} to produce a new
seed masked to 63 bits (\breaktt{\&\ 0x7FFFFFFFFFFFFFFF}, keeping it a
non-negative Python \texttt{int}), and
\breaktt{sample(grammar,\ count,\ base\_seed=None)} calls \texttt{expand}
once per derived seed to produce \texttt{count} independent
\texttt{Spec} objects --- independent in the sense that each is keyed
off a distinct derived seed, not off any shared mutable state.
\breaktt{enumerate\_all(grammar)} takes the orthogonal path: rather than
sampling, it walks the full \breaktt{itertools.product} over every slot's
options (reserved slots included) and returns one plain dict per cell,
giving direct access to the entire nominal product space for exhaustive
audits.
\end{block}

\begin{block}{Stage C --- Materialization into a runnable child tree}
\protect\phantomsection\label{stage-c-materialization-into-a-runnable-child-tree}
\breaktt{materialize(spec,\ out\_root,\ template\_root,\ clean=False)}
turns a resolved \texttt{Spec} into an actual directory of files on
disk. The child's directory name is derived deterministically by
\breaktt{child\_name(spec)} as
\texttt{child\_\{primitive\_domain\}\_\{spec\_hash\}} --- so the name
itself encodes both the selected domain and a content-derived identity,
without any counter or timestamp. \texttt{\_build\_tree} then assembles
the file map that will be written:

\begin{itemize}
\tightlist
\item
  \textbf{Kernel primitives.} \breaktt{\_vendor\_kernel\_sources} copies
  \breaktt{src/primitives/\ base.py} and
  \texttt{src/primitives/\{domain\}.py} for the spec's
  \breaktt{primitive\_domain}, rewriting \breaktt{from\ src.primitives} /
  \breaktt{from\ .primitives} imports to a bare \breaktt{from\ primitives}
  so the copied module resolves correctly once it is no longer nested
  under the parent template's package. It also synthesizes a minimal
  \breaktt{primitives/\_\_init\_\_.py} whose
  \breaktt{collect\_primitives()} imports only the one selected domain
  submodule --- the child ships exactly one domain's kernel, not all
  five.
\item
  \textbf{Figures source}, vendored verbatim from
  \texttt{src/figures.py} when present.
\item
  \textbf{Dependency vendoring.} \texttt{\_resolve\_deps} reads
  \texttt{dep\_mode} out of the spec's selections (defaulting to
  \texttt{"vendor"}) and, for each name in \texttt{spec.deps}, resolves
  the corresponding path from \breaktt{\_VENDORABLE\_MODULE\_PATHS}
  relative to the repository root
  (\breaktt{template\_root.parent.parent.parent}), embedding the real
  infrastructure source when it exists on disk or writing an explicit
  \breaktt{"Source\ not\ found\ at\ build\ time."} stub when it does not
  --- a vendored dependency the pipeline could not actually locate is
  marked as such in the generated file, not silently omitted. A seam
  \breaktt{vendored/\_\_init\_\_.py} re-exports every vendored module by
  name.
\item
  \textbf{A minimal \texttt{pyproject.toml}} declaring the child as its
  own installable project (\texttt{name\ =\ "child\_\{domain\}"},
  \texttt{numpy}/\texttt{matplotlib}/\texttt{pyyaml} runtime deps,
  pytest configured with \texttt{pythonpath\ =\ {[}"."{]}}).
\item
  \textbf{A generated \texttt{analysis.py}} that calls
  \breaktt{collect\_primitives()}, iterates the selected domain's
  registered \texttt{PrimitiveSpec} entries, and prints each one's name
  and result on its \texttt{example\_input} --- a real, executable entry
  point rather than a stub that only imports.
\item
  \textbf{A generated smoke test}, \breaktt{tests/test\_analysis.py},
  asserting only that \texttt{run()} completes without raising.
\item
  \textbf{Manuscript stubs} (\breaktt{\_emit\_manuscript}) --- abstract,
  introduction, results, and limitations sections written directly from
  the spec's own fields (domain, spec hash, all non-domain selections),
  so the child's own manuscript is itself traceable to the same
  \texttt{Spec} that generated its code.
\end{itemize}

Every file in the assembled tree is written to \texttt{child\_root} and
also retained in-memory as \texttt{written:\ dict{[}str,\ str{]}}.
\texttt{materialize} then calls
\breaktt{tree\_hash\_from\_content\_hashes(written)} --- sorting all
relative paths lexicographically, joining each as
\texttt{"\{path\}:\{content\}"}, and taking the SHA-256 of the
concatenation --- so the tree hash is a function of path names and byte
content only, not file-system metadata such as mtimes or insertion
order. The resulting \texttt{provenance.json} records a schema version
(\breaktt{"autopoiesis/provenance/1"}), the full
\texttt{spec.to\_dict()}, the computed tree hash, and the sorted list of
every written relative path --- the single artifact Stage D re-derives
from.
\end{block}

\begin{block}{Stage D --- Verification}
\protect\phantomsection\label{stage-d-verification}
\breaktt{verify\_child(child\_root)} (implemented in
\texttt{src/verify.py}, not \texttt{materialize.py}) loads
\texttt{provenance.json}, reads back every file named in its
\texttt{files} list, recomputes the tree hash from those live contents
via the same \breaktt{tree\_hash\_from\_content\_hashes} function used at
materialization time, and compares it against the recorded hash. Each
individual predicate --- \breaktt{provenance\_exists},
\breaktt{provenance\_parseable}, \breaktt{all\_files\_present},
\breaktt{tree\_hash\_matches} --- is captured as a
\breaktt{CheckResult(name,\ passed,\ detail)} inside an aggregate
\texttt{CheckReport}, so a caller can distinguish \emph{which} invariant
broke rather than receiving a single boolean. Any edit, deletion, or
addition to a file listed in provenance --- anything from a hand-edited
line to a regenerated timestamp --- changes at least one entry in
\texttt{live}, which changes the recomputed hash, which fails
\breaktt{tree\_hash\_matches} without requiring the verifier to inspect a
diff. \breaktt{verify\_child\_full} extends this with a
\breaktt{schema\_version\_correct} check against
\breaktt{materialize.PROVENANCE\_SCHEMA\_VERSION}, catching provenance
written by a different schema generation of the tool.
\end{block}

\begin{block}{Stage E --- Sealing}
\protect\phantomsection\label{stage-e-sealing}
\texttt{sealing.py} provides the payload and image layer for embedding a
tamper-evident summary alongside a generated child.
\breaktt{build\_payload(spec\_hash,\ tree\_hash,\ seed)} serializes a
compact JSON object binding all three identifiers;
\breaktt{build\_pointer\_payload} and \breaktt{build\_barcode\_payload}
offer lighter-weight alternatives (a URL-plus-hash pointer, and a
colon-joined label/hash/seed string, respectively) for contexts where a
full JSON payload is unnecessary. \texttt{qr\_matrix()} and
\texttt{qr\_image()} wrap the optional \texttt{qrcode} library, falling
back to a deterministic 5×5 checkerboard stub when it is absent, so the
sealing stage --- and its tests --- do not require an optional
dependency to be installed. \breaktt{verify\_seal(child\_root)} checks
for a \texttt{seal.json}, that it parses, and that it carries a
\texttt{spec\_hash} field, mirroring the same
\texttt{CheckResult}/\texttt{CheckReport} pattern used in Stage D rather
than introducing a parallel reporting shape.
\end{block}

\begin{block}{Property-based invariants}
\protect\phantomsection\label{property-based-invariants}
Beyond the fixed example-based tests enumerated in the Honesty
Contract's ground-truth table (§4),
\breaktt{tests/test\_property\_invariants.py} exercises the expansion and
materialization functions against Hypothesis-generated inputs using the
property-based testing paradigm \breakseq{[@claessen2000quickcheck;
@maciver2019hypothesis]}: rather than asserting fixed input/output
pairs, these tests assert invariants that must hold across a swept range
of seeds --- \texttt{product\_size} equals the literal product of every
slot's option count regardless of seed;
\breaktt{effective\_product\_size} does not exceed
\texttt{product\_size}; two grammars parsed from the same block produce
the same \texttt{grammar\_hash}; \texttt{expand()} called twice with the
same seed produces the same \texttt{spec\_hash} for any seed Hypothesis
draws in \texttt{{[}0,\ 10**9{]}}; the resolved
\breaktt{primitive\_domain} is one of \texttt{KNOWN\_DOMAINS}; two
independent \texttt{materialize()} calls from the same spec into
different output roots produce identical \texttt{tree\_hash} values, per
domain; and \texttt{verify\_child} reports \texttt{all\_passed} on a
freshly materialized child but reports a failure the moment any listed
file is edited or deleted, per domain. A parallel Hypothesis sweep over
generated text confirms \texttt{qr\_matrix()} returns a square matrix
and is deterministic on repeated calls with the same input.
\texttt{hypothesis} is declared under this project's \texttt{dev}
optional-dependency group in \texttt{pyproject.toml}, alongside
\texttt{pytest} and \texttt{pytest-cov} --- it is imported
unconditionally at the top of \breaktt{test\_property\_invariants.py}, so
it is a real (if dev-only) dependency of this test module, not a soft,
try/except guarded one.
\end{block}

\begin{block}{Reproducibility framing}
\protect\phantomsection\label{reproducibility-framing}
The tree-hash construction used at both Stage C and Stage D deliberately
mirrors the general goal of reproducible software builds
\breakseq{[@reproducible\_builds]}: a build (here, a materialize call) is
reproducible exactly when independent re-derivations from the same
inputs --- grammar, seed, template source --- produce byte-identical
output, and that identity is checked by content hash rather than by
trusting the process that produced it. The tree hash's construction ---
sort every \breaktt{(path,\ content)} pair lexicographically, join as
\texttt{"path:content"}, and take one SHA-256 of the concatenation ---
is a flat, single-level structure, not a binary hash tree.
\breaktt{src/integrity.py} separately exposes a genuine binary
\texttt{merkle\_root()} (pairwise concatenate-and-hash up the tree,
duplicating the final odd node when a level has odd cardinality), in the
spirit of the hash-tree provenance idea introduced for digital
signatures \breakseq{[@merkle\_tree\_provenance]}. As of this writing
\texttt{merkle\_root()} is a standalone, independently-tested utility:
\breaktt{integrity\_profile} is declared as a reserved grammar slot
(options \texttt{sha256}, \texttt{merkle} in this project's own
\texttt{config.yaml}; \texttt{SYNTAX.md} documents a longer illustrative
option list including \texttt{merkle\_kmyth}), but neither
\texttt{materialize()} nor \texttt{verify\_child()} currently branches
on its selected value --- the slot is present in the grammar and
excluded from the effective product size, but is not yet wired to switch
which hashing function actually runs. Reporting that precisely, rather
than implying the profile already gates behavior, is the same honesty
discipline the project asks of its own manuscript elsewhere.
\end{block}
\end{frame}

\end{document}
